\documentclass[10pt,a4paper]{article}
\usepackage[margin=0.75in]{geometry}
\usepackage{hyperref}
\usepackage{enumitem}

\title{\vspace{-1.5em}Digital Signal Encoder and Decoder}
\author{ITT 036 -- Programming Assignment 1}
\date{Autumn 2026}

\begin{document}
\maketitle
\vspace{-2em}

\section*{Overview}
This project is a C++ program that shows how digital signals are encoded
and how an analog signal is converted into digital form. The user can
either generate digital data or start from an analog signal, and the
program then draws the resulting waveforms on screen using Raylib.

\section*{Implementation}
On the digital side, the program can produce random binary data, data
with long runs of consecutive zeros, or data typed in by the user. It
also searches the generated sequence for its longest palindromic
substring. Five line coding schemes are implemented: NRZ-L, NRZ-I,
Manchester, Differential Manchester, and AMI. For AMI, the B8ZS and
HDB3 scrambling techniques are also available.

On the analog side, the program generates a sine wave and samples it.
Each sample is then quantized to one of eight levels. During PCM
encoding, every sample is turned into a 3-bit code, and the decoder
reverses this by rebuilding the quantized values from those codes.
Raylib is used throughout to plot both the digital signals and the
decoded output.

\section*{Tools and Assumptions}
The project is written in C++ and relies only on the C++ Standard
Library and Raylib. The analog signal is kept roughly between $-1$ and
$+1$. In its current form, the PCM stage takes 20 samples and uses
eight quantization levels, which gives a 60-bit encoded stream. Since
the decoder outputs the midpoint of each quantization level, the decoded
values will not always match the original samples exactly.

\section*{How to Run}
On Windows, you can run \texttt{assignment.exe} directly from the project
repository. Otherwise, \texttt{main.cpp} can be compiled with any C++
compiler that has been set up with Raylib. When the program starts, a
menu lets you choose digital input, analog input, or exit.

\section*{References and Acknowledgment}
\begin{enumerate}[leftmargin=*,nosep]
    \item ITT 036, Programming Assignment 1, Autumn 2026.
    \item Project source:
    \url{https://github.com/muntasir72/Data-Com-Assignment}
    \item Raylib documentation: \url{https://www.raylib.com/}
\end{enumerate}

I used AI assistance to help understand parts of the code and to
prepare this report.

\end{document}